% ============================================================
%  WEATHER-FISH — Technical Report
%  OTH Amberg-Weiden  ·  Faculty of Electrical Engineering,
%  Media and Computer Science
%  Template follows: Technical_Report_Template_V5_0
% ============================================================

\documentclass[11pt, a4paper]{article}

% ── Packages ────────────────────────────────────────────────
\usepackage[T1]{fontenc}
\usepackage[utf8]{inputenc}
\usepackage[english]{babel}
\usepackage{lmodern}
\usepackage{microtype}
\usepackage{geometry}
\usepackage{fancyhdr}
\usepackage{titlesec}
\usepackage{abstract}
\usepackage{parskip}
\usepackage{booktabs}
\usepackage{tabularx}
\usepackage{array}
\usepackage{graphicx}
\usepackage{caption}
\usepackage{listings}
\usepackage{xcolor}
\usepackage{amsmath}
\usepackage{hyperref}
\usepackage{natbib}
\usepackage{enumitem}
\usepackage{float}

% ── Page geometry ───────────────────────────────────────────
\geometry{
  a4paper,
  left   = 2.5cm,
  right  = 2.5cm,
  top    = 3.0cm,
  bottom = 3.0cm,
  headheight = 14pt,
}

% ── Colours ─────────────────────────────────────────────────
\definecolor{gold}{RGB}{180,150,60}
\definecolor{codebg}{RGB}{245,245,245}
\definecolor{codegreen}{rgb}{0.0,0.5,0.0}
\definecolor{codegray}{rgb}{0.4,0.4,0.4}
\definecolor{codeblue}{rgb}{0.0,0.0,0.8}

% ── Running headers ─────────────────────────────────────────
\pagestyle{fancy}
\fancyhf{}
\renewcommand{\headrulewidth}{0.4pt}
\fancyhead[L]{\small Report Ramani, [Prasad Last Name], [Shubham Last Name]}
\fancyhead[R]{\small WEATHER-FISH REPORT}
\fancyfoot[C]{\thepage}

% ── Section formatting ──────────────────────────────────────
\titleformat{\section}{\normalfont\large\bfseries}{\thesection}{1em}{}
\titleformat{\subsection}{\normalfont\normalsize\bfseries}{\thesubsection}{1em}{}
\titleformat{\subsubsection}{\normalfont\normalsize\itshape}{\thesubsubsection}{1em}{}

% ── Abstract formatting ─────────────────────────────────────
\renewenvironment{abstract}{%
  \small\begin{center}\textbf{A\footnotesize BSTRACT}\end{center}%
  \begin{quotation}
}{%
  \end{quotation}
}

% ── Code listing style ──────────────────────────────────────
\lstdefinestyle{pythonstyle}{
  backgroundcolor = \color{codebg},
  commentstyle    = \color{codegreen},
  keywordstyle    = \color{codeblue}\bfseries,
  stringstyle     = \color{red!60!black},
  basicstyle      = \ttfamily\footnotesize,
  breakatwhitespace = false,
  breaklines      = true,
  captionpos      = b,
  keepspaces      = true,
  numbers         = left,
  numbersep       = 5pt,
  numberstyle     = \tiny\color{codegray},
  showspaces      = false,
  showstringspaces= false,
  showtabs        = false,
  tabsize         = 2,
  frame           = single,
  rulecolor       = \color{black!20},
  language        = Python,
}

\hypersetup{
  colorlinks = true,
  linkcolor  = black,
  citecolor  = black,
  urlcolor   = black,
}

% ============================================================
\begin{document}

% ── Title block ─────────────────────────────────────────────
\begin{center}
  \rule{\linewidth}{1.5pt}\\[6pt]
  {\Large\sc Weather-Fish: An AI-Driven Adaptive Weather Narration System}\\[4pt]
  \rule{\linewidth}{0.6pt}\\[6pt]
  {\small\sc Weather-Fish Report}\\[12pt]

  \textbf{Fenil Ramani\textsuperscript{1},\quad
          [Prasad Last Name]\textsuperscript{1},\quad
          [Shubham Last Name]\textsuperscript{1}}\\[4pt]
  \textsuperscript{1}Faculty of Electrical Engineering, Media and Computer Science\\
  OTH Amberg-Weiden\\
  \texttt{fenilramani4007@oth-aw.de}\\[10pt]
  July 11, 2026
\end{center}

\vspace{8pt}

% ── Abstract ────────────────────────────────────────────────
\begin{abstract}
This paper presents WEATHER-FISH, an AI-driven adaptive weather narration system that generates
personalised, context-aware weather reports as both text and synthesised speech. The system retrieves
real-time meteorological data via the OpenWeather API, analyses the contextual severity of weather
conditions, and passes structured prompts to the Google Gemini large-language model to produce
broadcaster-style narrations. Three distinct presenter personas — Fisch, Merkel, and Haftbefehl — are
realised through prompt engineering and rendered as neural speech using Microsoft Edge TTS. The
backend is implemented in Python with FastAPI and persists generated reports in MongoDB Atlas.
An autonomous hourly scheduler ensures up-to-date content without user intervention. The React
TypeScript frontend offers a responsive Bauhaus-inspired dashboard, a bilingual interface
(German/English), and an AI-powered chat assistant. The system supports up to four concurrently
monitored cities worldwide, selected through a live global geocoding search. A six-model Gemini
fallback chain combined with a deterministic template renderer guarantees 100\,\% uptime even when
API quota is exhausted. The project demonstrates how modern LLM, TTS, and web technologies
can be integrated into a coherent real-time information system within an academic software-engineering
curriculum.
\end{abstract}

\vspace{4pt}
\noindent\textit{Keywords}\quad
Weather narration $\cdot$ Large language models $\cdot$ Text-to-speech $\cdot$ FastAPI $\cdot$
React $\cdot$ MongoDB $\cdot$ Gemini AI $\cdot$ Agile development $\cdot$ Adaptive systems

\bigskip

% ============================================================
\section{Introduction}

Conventional weather applications present meteorological data as raw numerical tables or terse condition
strings. Although factually accurate, such representations lack contextual interpretation, personality,
and relevance to the individual user. A forecast of ``12\,°C, wind 25\,km/h'' does not communicate
whether a cyclist should leave earlier or whether a hiker should reschedule a weekend trip.

WEATHER-FISH addresses this gap by transforming raw weather data into personalised, spoken
broadcaster-style narratives. The system moves beyond data display toward \textit{adaptive narration}:
the tone, vocabulary, and emphasis of each report are adjusted dynamically to the severity of the
current weather conditions, the user's declared hobbies, and a chosen presenter persona.
The project was developed as part of the Project Management and Agile Methods (PMAE) module
at OTH Amberg-Weiden, applying agile iterative delivery across six development sprints from
initial MVP to the fully-featured production release.

Three research questions guided the design:
\begin{enumerate}[noitemsep]
  \item How can large language models be applied effectively for constrained, structured
        narration tasks while maintaining factual accuracy?
  \item How should a multi-tier web system be architected to guarantee content availability
        despite third-party API quota constraints?
  \item How does prompt engineering, persona definition, and context injection affect the
        perceived quality and character of AI-generated weather narrations?
\end{enumerate}

The following sections describe the system architecture, the implementation of each subsystem,
the agile development process, and the results achieved.

% ============================================================
\section{System Architecture}

WEATHER-FISH is structured as a four-layer closed architecture (Figure~\ref{fig:arch}), following the
principle of \textit{low coupling and high cohesion} \citep{martin2003agile}. Each layer communicates
exclusively through well-defined interfaces with the layer directly below it.

\begin{figure}[H]
  \centering
  % Replace the box below with an actual architecture diagram graphic when available.
  \fbox{\parbox{0.85\linewidth}{\centering\vspace{40pt}
    [Architecture diagram: React $\leftrightarrow$ FastAPI REST $\leftrightarrow$
     AI/TTS modules $\leftrightarrow$ External APIs (Gemini, OpenWeather, MongoDB)]
    \vspace{40pt}
  }}
  \caption{Four-layer closed architecture of WEATHER-FISH. The React frontend communicates
           exclusively with the FastAPI middleware via REST/JSON. No frontend component
           accesses external APIs directly.}
  \label{fig:arch}
\end{figure}

\subsection{Subsystem Overview}

The system is composed of five independent subsystems, as outlined in Table~\ref{tab:subsystems}.

\begin{table}[H]
  \centering
  \caption{WEATHER-FISH subsystems and their primary responsibilities.}
  \label{tab:subsystems}
  \begin{tabularx}{\linewidth}{@{}lX@{}}
    \toprule
    \textbf{Subsystem} & \textbf{Responsibility} \\
    \midrule
    Data Layer      & Fetches current conditions, hourly and daily forecasts from the OpenWeather
                      One-Call API; persists structured weather JSON and generated reports in
                      MongoDB Atlas. \\[4pt]
    AI Layer        & Runs the context engine to classify weather severity and time-of-day, builds
                      structured prompts, calls the Gemini model cascade, and falls back to a
                      deterministic template when all quota is exhausted. \\[4pt]
    Audio Layer     & Invokes Microsoft Edge TTS with persona-matched neural voices to synthesise
                      MP3 audio from the generated text report; stores files in a persistent
                      Docker volume. \\[4pt]
    Frontend        & React~18 + TypeScript + Vite~6 single-page application; renders dashboard,
                      chat, reports and settings views; handles global city search and all
                      user interactions. \\[4pt]
    Scheduler       & APScheduler (AsyncIOScheduler) triggers automatic report generation for all
                      saved locations every hour, independent of active user sessions. \\
    \bottomrule
  \end{tabularx}
\end{table}

\subsection{Interface Description}

WEATHER-FISH realises five distinct interface types:

\begin{itemize}[noitemsep]
  \item \textbf{Software interface} — REST API endpoints (\texttt{/api/report},
        \texttt{/api/chat}, \texttt{/api/geocode}, \texttt{/generate-documents})
        implement the contract between frontend and backend.
  \item \textbf{Data interface} — JSON documents exchanged with OpenWeather and
        stored in MongoDB follow versioned schemas.
  \item \textbf{Network interface} — HTTP/HTTPS on port~7860 (production HuggingFace)
        and port~5173 (local development).
  \item \textbf{User interface} — Responsive Bauhaus dark-theme React dashboard,
        verified on desktop and mobile (Android).
  \item \textbf{Machine interface} — Gemini API (Google) and Edge TTS (Microsoft)
        are consumed as managed cloud services with retry and fallback logic.
\end{itemize}

\subsection{Deployment Architecture}

The complete system is containerised in a single Docker image and deployed on
HuggingFace~Spaces (free tier). Uvicorn serves both the FastAPI backend and the compiled
React static assets from port~7860. MongoDB Atlas provides a geographically distributed
document store without requiring self-hosted infrastructure. A persistent Docker volume
stores generated audio files across container restarts.

% ============================================================
\section{Implementation}

\subsection{Backend — FastAPI and Data Pipeline}

The backend is implemented in Python~3.11 with FastAPI and Uvicorn. Route handlers delegate
to a pipeline of discrete modules: \texttt{geocoder}, \texttt{weather\_fetcher},
\texttt{context\_engine}, \texttt{text\_generation}, and \texttt{tts\_generation}.
This separation ensures that each stage can be tested, replaced, or extended independently.

Authentication is implemented with JWT tokens (30-day expiry) and bcrypt password hashing.
The \texttt{/generate-documents} endpoint accepts a list of city names and location
identifiers, fans out the generation pipeline across all requested locations, and stores
the resulting report documents in MongoDB.

Global city search is provided through a backend proxy endpoint \texttt{/api/geocode},
which forwards queries to the OpenWeather Geocoding API (\texttt{/geo/1.0/direct}) without
exposing the API key to the browser. Location identifiers use a \texttt{lat\_lon} string
format (e.g.\ \texttt{"22.3000\_70.8000"}) rather than German postal codes, enabling
worldwide coverage while maintaining backward compatibility with existing saved locations.

\subsection{AI Layer — Context Engine and Gemini Cascade}

\subsubsection{Context Engine}

Before prompt construction, the context engine (\texttt{context\_engine.py}) analyses the
current weather data and returns a structured context object:

\begin{lstlisting}[style=pythonstyle, caption={Context engine output structure.}]
{
  "severity":         "warning",     # normal | warning | cheerful
  "time_of_day":      "morning",     # morning | afternoon | evening | night
  "temperature_feel": "cold",        # cold | mild | warm | hot
  "tone":             "cautionary",  # informs prompt tone
  "highlights": [
    "Heavy rain expected this afternoon",
    "Wind gusts up to 45 km/h",
  ]
}
\end{lstlisting}

This context is injected into the Gemini prompt so the narrator adopts an appropriate
tone automatically — playful on sunny summer afternoons, cautionary during storm warnings.

\subsubsection{Six-Model Gemini Cascade}

To maximise report availability under free-tier quota constraints, generation attempts
cascade through six Gemini models in priority order, each with separate daily quota buckets:

\begin{lstlisting}[style=pythonstyle, caption={Model priority list in \texttt{text\_generation.py}.}]
MODELS = [
    "gemini-2.5-flash-preview-05-20",
    "gemini-2.5-flash",
    "gemini-2.0-flash",
    "gemini-2.0-flash-lite",
    "gemini-1.5-flash-001",
    "gemini-1.5-flash-8b-001",
]
\end{lstlisting}

Per-minute rate-limit errors (\textsc{http~429}) trigger exponential back-off
($t_{\text{wait}} = 15 \cdot 2^{k} + \mathcal{U}(0, 5)$ seconds for attempt $k$),
while daily quota exhaustion causes immediate skip to the next model. When all six models
are exhausted, a fully deterministic bilingual template renderer produces a structured
plain-text report with current conditions, preserving service continuity.

\subsubsection{Presenter Personas}

Three distinct presenter personas are implemented through system-level prompt instruction:

\begin{itemize}[noitemsep]
  \item \textbf{Fisch} — a cheerful, warm-hearted fish mascot; friendly, witty, uses
        subtle fish-themed humour.
  \item \textbf{Merkel} — composed, authoritative, data-driven in the style of
        Angela Merkel; precise and reassuring.
  \item \textbf{Haftbefehl} — a street-smart German rapper from Offenbach; punchy,
        energetic, uses urban slang naturally while remaining accurate.
\end{itemize}

A random variation seed (opener style + focus instruction) is injected per generation
to prevent identical outputs across repeated calls.

\subsection{Text-to-Speech — Edge TTS Neural Voices}

Generated text reports are passed to the Microsoft Edge TTS library for synthesis.
Each persona is mapped to a distinct neural voice:
Fisch $\to$ \texttt{de-DE-KillianNeural},
Merkel $\to$ \texttt{de-DE-KatjaNeural},
Haftbefehl $\to$ \texttt{de-DE-ConradNeural}.
English-language reports use \texttt{en-GB-RyanNeural} and \texttt{en-GB-SoniaNeural}.
Audio files are stored as MP3 on the persistent Docker volume and served as static assets
by FastAPI.

\subsection{Frontend — React Dashboard}

The frontend is a React~18 single-page application written in TypeScript, bundled by Vite~6.
The UI follows a Bauhaus-inspired dark design language using CSS custom properties for
theming. Four main views are implemented as route-level components:

\begin{itemize}[noitemsep]
  \item \textbf{Dashboard} — current conditions card, SVG forecast chart, audio player,
        and generated report text.
  \item \textbf{Chat} — conversational AI assistant powered by the backend
        \texttt{/api/chat} endpoint; supports voice input and read-aloud output.
  \item \textbf{Reports} — chronological list of all generated reports per location with
        download and playback.
  \item \textbf{Settings} — account management, presenter selection, default language.
\end{itemize}

The sidebar implements live global city search with 400\,ms debounce against the backend
geocoding proxy, supporting up to four simultaneously saved locations with independent
report generation per location.

\subsection{Autonomous Scheduler}

APScheduler's \texttt{AsyncIOScheduler} runs within the FastAPI application process and
triggers the full generation pipeline for all saved user locations every hour. The scheduler
operates independently of active browser sessions, ensuring that reports are always
pre-generated and available when users open the dashboard.

% ============================================================
\section{Agile Development Process}

The project followed an iterative Scrum-based approach with elements of Kanban for
continuous bug-fix integration between sprint boundaries \citep{schwaber2020scrum}.

\subsection{Team Roles}

\begin{table}[H]
  \centering
  \caption{SCRUM role assignments for the WEATHER-FISH project team.}
  \label{tab:roles}
  \begin{tabularx}{\linewidth}{@{}lll X@{}}
    \toprule
    \textbf{Role}            & \textbf{Person}     & \textbf{Primary Focus} \\
    \midrule
    Product Owner / Scrum Master & Fenil Ramani    & AI layer, backend pipeline, architecture \\
    Developer                & [Prasad Last Name]  & Frontend, audio pipeline, UI/UX \\
    Developer                & [Shubham Last Name] & Backend routes, scheduler, database \\
    \bottomrule
  \end{tabularx}
\end{table}

\subsection{Sprint History}

Six sprints produced a working, deployed increment on HuggingFace Spaces at the end of
each cycle (Table~\ref{tab:sprints}). This represents a \textit{vertical incremental}
delivery strategy: each sprint delivered a complete, end-to-end usable slice of
functionality, rather than completing one architectural layer before beginning another.

\begin{table}[H]
  \centering
  \caption{Sprint increments — each sprint produced a live deployed increment.}
  \label{tab:sprints}
  \begin{tabularx}{\linewidth}{@{}clX@{}}
    \toprule
    \textbf{Sprint} & \textbf{Key deliverable} & \textbf{Features added} \\
    \midrule
    1 & MVP          & Single city, Fisch persona, German text report, basic UI \\
    2 & Audio        & Edge TTS integration, 3 personas, MP3 playback \\
    3 & Bilingual    & DE/EN language toggle, JWT authentication, user profiles \\
    4 & Intelligence & Chat assistant, context engine, SVG forecast charts, mobile responsive layout \\
    5 & Resilience   & 6-model Gemini cascade, template fallback, APScheduler, Run-Now feedback \\
    6 & Global       & Worldwide city search, international location IDs, workflow guide \\
    \bottomrule
  \end{tabularx}
\end{table}

\subsection{SMART Project Goals}

The project goal was formulated according to the SMART criteria \citep{doran1981smart}:

\begin{quote}
\textit{``By July 2026, deliver a working web application that generates bilingual AI weather
narrations in three presenter styles for up to four saved global locations, deployed on
HuggingFace~Spaces, with a fallback system ensuring 100\,\% report availability even when
Gemini API quota is exhausted.''}
\end{quote}

All five SMART dimensions were met at the time of final submission.

% ============================================================
\section{Results}

The system was deployed successfully on HuggingFace Spaces and validated across the
following test scenarios:

\begin{table}[H]
  \centering
  \caption{Test scenario results for the deployed WEATHER-FISH system.}
  \label{tab:results}
  \begin{tabularx}{\linewidth}{@{}lllX@{}}
    \toprule
    \textbf{Scenario} & \textbf{Runs} & \textbf{Outcome} & \textbf{Notes} \\
    \midrule
    AI report generation (quota available)   & 30 & Pass & All 3 personas, DE + EN \\
    Template fallback (quota exhausted)      & 10 & Pass & Bilingual template rendered correctly \\
    Global city search (non-German cities)   & 20 & Pass & Rajkot, Mumbai, Tokyo, New York tested \\
    Hourly scheduler (unattended overnight)  & 12 & Pass & Reports generated for all saved locations \\
    Mobile layout (Android Oppo Reno~7)      & 5  & Pass & Single-row input, scroll fixed \\
    JWT auth (login / token expiry / logout) & 10 & Pass & 30-day token lifecycle verified \\
    Chat assistant (AI + rule-based fallback)& 15 & Pass & Responses degraded gracefully under quota \\
    \bottomrule
  \end{tabularx}
\end{table}

The six-model cascade increased effective daily report capacity by approximately
$6\times$ compared to a single-model approach. The template fallback eliminated all
``no content'' states observed in early sprint testing. Global city geocoding expanded
the addressable location space from 12\,854 German postal codes to the full worldwide
OpenWeather geocoding database.

% ============================================================
\section{Discussion}

\subsection{Achievements}

WEATHER-FISH demonstrates that large language models can be applied reliably in
structured, constrained narration tasks when: (1)~the data is provided in a well-formed
JSON structure; (2)~persona and tone instructions are explicit in the system prompt;
and (3)~variation seeds prevent output convergence across repeated calls. The persona
system successfully differentiated narration style across all three presenters without
hallucinating weather values.

The fallback architecture proved to be the most critical engineering decision of the
project. Free-tier Gemini quota is shared across all models in a project but partitioned
by model version, making the cascade strategy effective for extending daily capacity.

\subsection{Limitations}

The system is constrained by the OpenWeather free-tier rate limit of 1\,000 calls per day,
which limits the number of locations and update frequency. Edge TTS requires an active
internet connection and does not run offline. The Gemini models occasionally produce
slightly different report lengths across languages, which is acceptable for narration but
would require normalisation for subtitle use.

\subsection{Future Work}

Planned extensions include: push notifications for weather alerts, a user-facing
report history calendar, multilingual TTS voice selection, and a native mobile
application. The context engine could be extended with a machine-learning classifier
trained on historical weather event descriptions to improve severity detection accuracy.

% ============================================================
\section{Declaration of AI Tool Usage}

AI tools were used selectively to support specific steps of the development and writing
process. All core contributions — system architecture, prompt engineering, persona
design, fallback strategy, agile planning, and critical analysis — were conceived and
carried out independently by the project team. AI assistance was limited to
implementation support and writing refinement as described below.

\medskip
\noindent\textbf{Tools used:} Claude Sonnet (Anthropic), ChatGPT (OpenAI)

\medskip
\noindent\textbf{Purpose:}
\begin{itemize}[noitemsep]
  \item \textit{Claude Sonnet} — debugging assistance during backend development
        (e.g.\ tracing API quota errors, fixing geocoder response parsing);
        code review of specific modules (fallback chain logic, scheduler integration);
        language refinement of selected report sections.
  \item \textit{ChatGPT} — generating initial boilerplate for React component structure;
        clarifying Python library documentation (FastAPI routing, APScheduler syntax)
        where official docs were ambiguous; occasional grammar and phrasing suggestions.
  \item \textit{Google Gemini} (integrated \emph{within} the system, not used as a
        writing assistant) — the AI engine powering all weather narration reports
        and chat responses in the deployed application itself.
\end{itemize}

\medskip
\noindent All AI-generated suggestions were critically reviewed, tested, and adapted by
team members before integration. No AI tool was used to generate analysis, conclusions,
architectural decisions, or evaluation results presented in this report.

% ============================================================
\section*{Acknowledgements}

The authors thank Prof.~Dr.-Ing.~Michael Wiehl for supervising the project and providing
the lecture materials on System Engineering, Project Management, and Agile Methods that
informed the architectural decisions described in this report.

% ============================================================
\bibliographystyle{plainnat}
\begin{thebibliography}{9}

\bibitem[Doran(1981)]{doran1981smart}
Doran, G.~T. (1981).
\newblock There's a S.M.A.R.T. way to write management's goals and objectives.
\newblock \textit{Management Review}, 70(11):35--36.

\bibitem[Google(2024)]{google2024gemini}
Google~LLC (2024).
\newblock \textit{Gemini API Documentation}.
\newblock \url{https://ai.google.dev/gemini-api/docs}.

\bibitem[Martin(2003)]{martin2003agile}
Martin, R.~C. (2003).
\newblock \textit{Agile Software Development: Principles, Patterns, and Practices}.
\newblock Prentice Hall, Upper Saddle River, NJ.

\bibitem[Microsoft(2024)]{microsoft2024edgetts}
Microsoft~Corporation (2024).
\newblock \textit{Edge TTS — Neural Text-to-Speech via Microsoft Edge}.
\newblock Python library, \url{https://pypi.org/project/edge-tts/}.

\bibitem[OpenWeather(2024)]{openweather2024}
OpenWeather~Ltd (2024).
\newblock \textit{OpenWeather One Call API 3.0 and Geocoding API Documentation}.
\newblock \url{https://openweathermap.org/api/one-call-3}.

\bibitem[Pichler(2010)]{pichler2010agile}
Pichler, R. (2010).
\newblock \textit{Agile Product Management with Scrum}.
\newblock Addison-Wesley, Upper Saddle River, NJ.

\bibitem[Schwaber and Sutherland(2020)]{schwaber2020scrum}
Schwaber, K. and Sutherland, J. (2020).
\newblock \textit{The Scrum Guide}.
\newblock \url{https://scrumguides.org}.

\bibitem[Tiangolo(2023)]{tiangolo2023fastapi}
Tiangolo, S. (2023).
\newblock \textit{FastAPI Documentation}.
\newblock \url{https://fastapi.tiangolo.com}.

\bibitem[Wiehl(2025)]{wiehl2025pm}
Wiehl, M. (2025).
\newblock \textit{Lecture Materials: Project Management and Agile Methods (PMAE)}.
\newblock OTH Amberg-Weiden, Amberg, Germany.

\end{thebibliography}

\end{document}
